\documentclass[11pt]{article}

\usepackage[margin=1in]{geometry}
\usepackage[T1]{fontenc}
\usepackage{lmodern}
\usepackage{microtype}
\usepackage{amsmath,amssymb,amsthm,mathtools,bm,mathrsfs}
\usepackage{centernot}
\usepackage{booktabs,array,tabularx}
\usepackage{enumitem}
\usepackage{xcolor}
\usepackage[hidelinks]{hyperref}
\usepackage[nameinlink,noabbrev]{cleveref}

\newtheorem{theorem}{Theorem}[section]
\newtheorem{proposition}[theorem]{Proposition}
\newtheorem{lemma}[theorem]{Lemma}
\newtheorem{corollary}[theorem]{Corollary}
\newtheorem{conjecture}[theorem]{Conjecture}
\theoremstyle{definition}
\newtheorem{definition}[theorem]{Definition}
\newtheorem{problem}[theorem]{Problem}
\theoremstyle{remark}
\newtheorem{remark}[theorem]{Remark}

\newcommand{\ket}[1]{\left|#1\right\rangle}
\newcommand{\bra}[1]{\left\langle#1\right|}
\newcommand{\I}{\mathrm i}
\newcommand{\e}{\mathrm e}
\newcommand{\Id}{\mathbb I}
\newcommand{\R}{\mathbb R}
\newcommand{\Z}{\mathbb Z}
\newcommand{\E}{\mathbb E}
\newcommand{\Prob}{\mathbb P}
\newcommand{\cH}{\mathcal H}
\newcommand{\cT}{\mathcal T}
\newcommand{\cL}{\mathcal L}
\newcommand{\cC}{\mathcal C}
\newcommand{\cK}{\mathcal K}
\newcommand{\cE}{\mathcal E}
\newcommand{\cA}{\mathcal A}
\newcommand{\ran}{\operatorname{ran}}
\newcommand{\Span}{\operatorname{span}}
\newcommand{\diag}{\operatorname{diag}}
\newcommand{\tr}{\operatorname{tr}}
\newcommand{\Var}{\operatorname{Var}}
\newcommand{\dist}{\operatorname{dist}}
\newcommand{\ip}[2]{\left\langle #1,#2\right\rangle}
\newcommand{\norm}[1]{\left\lVert #1\right\rVert}
\newcommand{\abs}[1]{\left\lvert #1\right\rvert}
\newcommand{\dd}{\,\mathrm d}
\newcommand{\eps}{\varepsilon}

\title{\textbf{Quenched Pair-Coupling Disorder in Frustration-Free QGN Parents II}\\[0.35em]
\large Quenched dynamical-source fluctuations, sharp cutoffs, and disorder-noise geometry}
\author{Nicholas Sledgianowski}
\date{Reviewer-integrated split draft --- July 27, 2026}

\begin{document}
\maketitle

\begin{abstract}
This companion paper studies quenched fluctuations of the low-energy dynamical source tail for the periodic random-conductance operator arising from the weighted QGN parent of Paper~I.  For product i.i.d. conductances on $\Lambda_L$, the physical sharp observable is
\[
 T_L(\delta)=L^{-d}\ip{g_L}{L_{J,L}^+\bm1_{(0,\delta]}(L_{J,L})g_L}.
\]
The spectral projector is discontinuous in each bond variable.  We therefore introduce the positive Abel envelope
$\mathfrak A_L(\delta)=L^{-d}\langle g_L,L_{J,L}^+e^{-L_{J,L}/\delta}g_L\rangle$ and derive an exact bond-sensitivity identity.  Under a one-bond Poincar\'e inequality, a period-uniform gradient-propagator estimate gives
\[
 \Var\mathfrak A_L(\delta)\le Cj_+^2\abs A^4L^{-d}(\delta/j_+)^{d/2}.
\]
This yields a one-sided quenched bound for $T_L$.  The same second moment gives the target variance scale at every bounded effective-mode cutoff only through $\Var T_L\le\E T_L^2$; it is not a mesoscopic concentration theorem.  In the frozen-spectrum Born sector, the sharp tail is an exact quadratic form and obeys Hanson--Wright concentration.

We audit three routes to a nonlinear fixed-cutoff theorem.  A finite discrete relative Wegner estimate controls eigenvalue counts but a modulated cycle shows that counts do not control current weight.  Exponential threshold randomization is exactly the Abel envelope, while positive heat-mixture thresholds cannot be relatively narrow.  Logarithmic shell averaging gives the optimal shell scale and a deterministic interval argument removes exceptional cutoffs only at the uncentered floor.  Rank-one interlacing proves that one bond crosses a prescribed cutoff at most once; Krein coordinates and a coarea identity determine the exact current-weighted crossing surface.  An acoustic surrogate shows why a vanishing uniform maximal-shell modulus is too strong.  The remaining mesoscopic theorem is isolated as a fixed-cutoff disorder-noise quadratic-variation bound with an integrable $t^{-1/2}$ boundary layer.
\end{abstract}

\paragraph{Claim hierarchy.}
\begin{enumerate}[leftmargin=2.2em,itemsep=0.2em]
 \item The exact QGN-to-random-conductance reduction and annealed source-tail estimate are stated as inputs from Paper~I.
 \item The Abel sensitivity identity, gradient-propagator reduction, smooth-tail variance, and one-sided sharp-tail theorem are proved here.
 \item The bounded-mode sharp variance uses only $\Var T\le\E T^2$ and is not claimed as mesoscopic centering.
 \item The Born concentration, relative Wegner estimate, randomization identities, log-shell theorem, one-crossing geometry, Krein formulas, participation coarea, and noise-semigroup identities are proved here.
 \item The remaining open input is the fixed-cutoff disorder-noise estimate stated in Problem~\ref{prob:fixed-cutoff-noise}; no forward equation reference is used in this front matter.
\end{enumerate}

\section{Setting and inputs from Paper I}
\label{sec:setting}

Let $d\ge2$, $\Lambda_L=(\Z/L\Z)^d$, $N=L^d$, and orient each nearest-neighbor bond in the positive coordinate direction.  The conductances on the $dL^d$ torus bonds are independent with common law $\beta\subset[j_-,j_+]$ and are extended periodically.  This is the product-torus ensemble, not a coupled periodization of one infinite realization.  One-dimensional rings appear below only as exactly solvable obstructions.

For $A\in\R^d$, define $a_{x,\mu}=A_\mu$,
\begin{align}
 g_L(x)&=\sum_{\mu=1}^d\bigl(j_{x-e_\mu,\mu}-j_{x,\mu}\bigr)A_\mu,\\
 (L_{J,L}f)(x)&=\sum_{\mu=1}^d\left[j_{x,\mu}(f(x)-f(x+e_\mu))+j_{x-e_\mu,\mu}(f(x)-f(x-e_\mu))\right].
\end{align}
Thus $g_L=BJ_LRA$ and $L_{J,L}=BJ_LB^*$.  Let $\tau_x$ translate the periodic environment and define the horizontal differences
\begin{equation}
 D_\mu F(J)=F(\tau_{e_\mu}J)-F(J),
 \qquad
 D_\mu^*F(J)=F(\tau_{-e_\mu}J)-F(J).
 \end{equation}
The environment generator and local drift are
\begin{equation}
 \mathscr L_L=D^*j(0)D,
 \qquad
 d_A=D^*(j(0)A).
 \label{eq:environment-operator-definition}
\end{equation}
The stationary extension of $d_A$ is $g_L$.  Paper~I proves the exact site/target identity and the total-energy relation
\begin{equation}
 \ip{g_L}{L_{J,L}^+g_L}=\sum_ej_ea_e^2-\cE_{J_L}[a].
 \label{eq:total-Hminus1}
\end{equation}
It also proves the deterministic gap
\begin{equation}
 \lambda_{\rm gap,L}=4j_-\sin^2(\pi/L)
 \label{eq:finite-size-gap}
\end{equation}
and the period-uniform annealed estimates
\begin{align}
 \E N^{-1}\ip{g_L}{\e^{-tL_{J,L}}g_L}
 &\le Cj_+^2\abs A^2(1+j_+t)^{-1-d/2},
 \label{eq:uniform-annealed-heat}\\
 \E T_L(\delta)&\le Cj_+\abs A^2\min\{1,(\delta/j_+)^{d/2}\}.
 \label{eq:uniform-finite-torus-tail}
\end{align}
We restate these inputs to make the fluctuation arguments self-contained at the operator level; the many-body derivation is in Paper~I \cite{Sledgianowski2026RobustnessI}.

\begin{theorem}[Imported exact and annealed input from Paper I]
\label{thm:quantitative-finite-torus}
Equations \eqref{eq:total-Hminus1}--\eqref{eq:uniform-finite-torus-tail} hold uniformly in $L$.  The exact environment--site intertwining identifies the stationary extension of $\e^{-t\mathscr L_L}d_A$ with $\e^{-tL_{J,L}}g_L$.
\end{theorem}
\section{Abel smoothing, variance, and the sharp edge}
\label{sec:quenched-fluctuations}

The annealed theorem controls the mean of the physical sharp tail
\begin{equation}
 T_L(\delta)
 :=\frac1{L^d}\ip{g_L}{L_{J,L}^+
 \bm{1}_{(0,\delta]}(L_{J,L})g_L}.
 \label{eq:sharp-tail-definition}
\end{equation}
A centered fluctuation proof cannot begin by differentiating this expression:
the map $J\mapsto\bm 1_{(0,\delta]}(L_{J,L})$ jumps when an eigenvalue crosses
the cutoff.  We first replace the projector by a positive smooth envelope, prove
the optimal effective-mode variance scale for that envelope, and then return to
the sharp observable by a samplewise comparison.

Define the Abel tail
\begin{align}
 \mathfrak A_L(\delta)
 &:=\frac{1}{L^d}\ip{g_L}{L_{J,L}^+
 \e^{-L_{J,L}/\delta}g_L}
 \label{eq:Abel-tail-definition}\\
 &= \frac{1}{L^d}\int_{1/\delta}^{\infty}
 \ip{g_L}{\e^{-tL_{J,L}}g_L}\dd t.
 \label{eq:Abel-tail-heat}
\end{align}
Since $\e^{-\lambda/\delta}\ge\e^{-1}$ for
$0<\lambda\le\delta$, one has the realization-wise envelope
\begin{equation}
 \boxed{T_L(\delta)\le \e\,\mathfrak A_L(\delta).}
 \label{eq:Abel-dominates-sharp}
\end{equation}
Integrating \eqref{eq:uniform-annealed-heat} also gives
\begin{equation}
 \E\mathfrak A_L(\delta)
 \le Cj_+\abs A^2
 \left(\frac{\delta}{j_+}\right)^{d/2},
 \qquad 0<\delta\le j_+.
 \label{eq:Abel-mean}
\end{equation}

The envelope has an exact cell-problem interpretation.  Put
\begin{equation}
 \cE_J(\phi)=\frac{1}{L^d}
 \sum_ej_e\bigl(a_e+(B^*\phi)_e\bigr)^2,
 \qquad
 \phi_*=-L_{J,L}^+g_L,
\end{equation}
and define the parabolic corrector
\begin{equation}
 \phi_t=-(\Id-\e^{-tL_{J,L}})L_{J,L}^+g_L.
\end{equation}

\begin{lemma}[Parabolic energy-excess identity]
\label{lem:parabolic-energy-excess}
For every realization and $t\ge0$,
\begin{equation}
 \boxed{
 \cE_J(\phi_t)-\cE_J(\phi_*)
 =\frac{1}{L^d}\ip{g_L}{L_{J,L}^+
 \e^{-2tL_{J,L}}g_L}.}
 \label{eq:parabolic-energy-excess}
\end{equation}
Hence $\mathfrak A_L(\delta)$ is exactly the residual cell-energy error after
parabolic relaxation for time $(2\delta)^{-1}$.
\end{lemma}

\begin{proof}
Completing the square around the minimizer gives
\[
 \cE_J(\phi)-\cE_J(\phi_*)
 =\frac{1}{L^d}\ip{\phi-\phi_*}{L_{J,L}(\phi-\phi_*)}.
\]
Here $\phi_t-\phi_*=\e^{-tL_{J,L}}L_{J,L}^+g_L$.  The spectral theorem yields
\eqref{eq:parabolic-energy-excess}.
\end{proof}

A rational positive envelope will also be useful.  Define
\begin{equation}
 \mathfrak R_{L,1}(\delta)
 =\frac{1}{L^d}\ip{g_L}
 {\frac{\delta^2}{L_{J,L}(L_{J,L}+\delta)^2}g_L}.
 \label{eq:Richardson-one-tail}
\end{equation}
Because $\delta^2/(\lambda+\delta)^2\ge1/4$ on
$0<\lambda\le\delta$,
\begin{equation}
 \boxed{T_L(\delta)\le4\mathfrak R_{L,1}(\delta).}
 \label{eq:Richardson-dominates-sharp}
\end{equation}
This is the first positive modified-corrector/Richardson spectral bias of
Gloria and Mourrat \cite{GloriaMourrat2012}.

\subsection{Exact vertical derivative and the gradient propagator}

Let $b_e=B\mathbf e_e$ be the incidence column of edge $e$ and set
\begin{equation}
 w_e(t)=b_e^*\e^{-tL_{J,L}}g_L,
 \qquad \tau=\delta^{-1}.
 \label{eq:edge-gradient-heat-source}
\end{equation}

\begin{lemma}[Exact Abel-tail sensitivity]
\label{lem:Abel-sensitivity}
At every coefficient configuration in the interior of the ellipticity box,
\begin{equation}
 \boxed{
 L^d\partial_{j_e}\mathfrak A_L(\delta)
 =2a_e\int_\tau^\infty w_e(s)\dd s
 -\int_\tau^\infty\int_0^s
 w_e(s-r)w_e(r)\dd r\dd s.}
 \label{eq:Abel-sensitivity}
\end{equation}
\end{lemma}

\begin{proof}
Use \eqref{eq:Abel-tail-heat},
$\partial_{j_e}g_L=a_eb_e$, and
$\partial_{j_e}L_{J,L}=b_eb_e^*$.  Duhamel's formula gives
\[
 \partial_{j_e}\e^{-sL_{J,L}}
 =-\int_0^s\e^{-(s-r)L_{J,L}}b_eb_e^*
 \e^{-rL_{J,L}}\dd r.
\]
Self-adjointness turns the two matrix elements in the second term into
$w_e(s-r)w_e(r)$.
\end{proof}

The new probabilistic input is a gradient upgrade of the divergence-semigroup
bound already used.  Related environment-process decay estimates are discussed in
\cite{deBuyerMourrat2015}; the precise period-uniform statements and the
weighted-Hölder argument needed here are collected in
Appendix~\ref{app:gradient-propagator}.

\begin{lemma}[Periodic gradient-propagator estimate]
\label{lem:periodic-gradient-propagator}
Under the periodic i.i.d. uniformly elliptic hypotheses, for $p=1,2$,
\begin{equation}
 \left(
 \E\frac{1}{dL^d}\sum_e\abs{w_e(t)}^{2p}
 \right)^{1/(2p)}
 \le C_pj_+\abs A(1+j_+t)^{-1-d/4}.
 \label{eq:periodic-gradient-propagator}
\end{equation}
The same conclusion holds for a periodic mixing ensemble whenever the
high-moment divergence-semigroup and weighted parabolic Green-gradient
estimates used below are uniform in the period.
\end{lemma}

\begin{proof}
Appendix~\ref{app:gradient-propagator} defines the torus Green function and weight, restates the two imported GNO estimates with their exact hypotheses and source numbers, and gives the weighted-H\"older calculation in full.  The constants are uniform in $L$.
\end{proof}

\subsection{CLT-order variance of the smooth tail}

For the centered theorem we impose one additional, explicit regularity
assumption on the one-edge law $\beta$.  Suppose
\begin{equation}
 \Var_\beta f
 \le C_\beta\int\abs{f'(s)}^2\dd\beta(s),
 \qquad C_\beta\le C_0j_+^2.
 \label{eq:one-edge-Poincare}
\end{equation}
The product law then tensorizes.  Uniform and smooth positive densities on a
fixed ellipticity interval satisfy such an inequality.  Discrete laws require
an oscillation/Efron--Stein version and are not included in the derivative
statement below.

\begin{theorem}[Quenched Abel-tail variance]
\label{thm:quenched-Abel-variance}
Assume the hypotheses of \cref{lem:periodic-gradient-propagator} and the
one-edge Poincar\'e inequality \eqref{eq:one-edge-Poincare}.  For
$0<\delta\le j_+$,
\begin{equation}
 \boxed{
 \Var\mathfrak A_L(\delta)
 \le Cj_+^2\abs A^4L^{-d}
 \left(\frac{\delta}{j_+}\right)^{d/2}.}
 \label{eq:quenched-Abel-variance}
\end{equation}
The constant depends only on $d$, $j_-/j_+$, and $C_0$.
\end{theorem}

\begin{proof}
Normalize $j_+=1$ and put
$h(t)=(1+t)^{-1-d/4}$.  Minkowski's inequality and
\eqref{eq:periodic-gradient-propagator} with $p=1$ give
\begin{equation}
 \left(
 \E\frac{1}{dL^d}\sum_e
 \left|\int_\tau^\infty w_e(s)\dd s\right|^2
 \right)^{1/2}
 \le C\abs A\int_\tau^\infty h(s)\dd s
 \le C\abs A\tau^{-d/4}.
 \label{eq:linear-sensitivity-tail}
\end{equation}
For the quadratic term, use $p=2$ and H\"older in probability:
\[
 \norm{w_e(r)w_e(u)}_2
 \le\norm{w_e(r)}_4\norm{w_e(u)}_4.
\]
Moreover,
\begin{equation}
 \int_{r,u\ge0:\ r+u\ge\tau}h(r)h(u)\dd r\dd u
 \le2\norm h_1\int_{\tau/2}^\infty h(s)\dd s
 \le C\tau^{-d/4}.
 \label{eq:convolution-tail}
\end{equation}
Applying Minkowski once more therefore bounds the averaged $L^2(\Omega_L)$
norm of the double integral in \eqref{eq:Abel-sensitivity} by
$C\abs A^2\tau^{-d/4}$.  Since $\abs{a_e}\le\abs A$,
\begin{equation}
 \frac{1}{dL^d}\sum_e
 \E\abs{\partial_{j_e}\mathfrak A_L(\delta)}^2
 \le CL^{-2d}\abs A^4\tau^{-d/2}.
 \label{eq:Abel-derivative-square}
\end{equation}
Tensorized Poincar\'e and summation over the $dL^d$ bonds give
$\Var\mathfrak A_L\le C\abs A^4L^{-d}\delta^{d/2}$.
Finally, $L_{J,L}=j_+L_{\widetilde J,L}$ and
$g_L=j_+g_{\widetilde J,L}$ imply
\[
 \mathfrak A_L^J(\delta)
 =j_+\mathfrak A_L^{\widetilde J}(\delta/j_+),
\]
which restores \eqref{eq:quenched-Abel-variance}.
\end{proof}

\begin{remark}[Higher moments]
Combining the same sensitivity estimate with a tensorized high-moment
Poincar\'e inequality gives, for every fixed $p<\infty$,
\begin{equation}
 \norm{\mathfrak A_L-\E\mathfrak A_L}_{L^{2p}}
 \le C_pj_+\abs A^2L^{-d/2}
 \left(\frac{\delta}{j_+}\right)^{d/4}.
 \label{eq:Abel-higher-moment}
\end{equation}
We do not claim an exponential tail without separately auditing the growth of
$C_p$.
\end{remark}

\begin{corollary}[Quenched one-sided sharp-tail bound]
\label{cor:quenched-sharp-tail}
Under the hypotheses of \cref{thm:quenched-Abel-variance}, for every
$0<\vartheta<1$, with probability at least $1-\vartheta$,
\begin{equation}
 \boxed{
 T_L(\delta)
 \le Cj_+\abs A^2\left[
 \left(\frac{\delta}{j_+}\right)^{d/2}
 +\vartheta^{-1/2}L^{-d/2}
 \left(\frac{\delta}{j_+}\right)^{d/4}
 \right].}
 \label{eq:quenched-sharp-tail}
\end{equation}
In addition,
\begin{equation}
 \E T_L(\delta)^2
 \le Cj_+^2\abs A^4\left[
 \left(\frac{\delta}{j_+}\right)^d
 +L^{-d}\left(\frac{\delta}{j_+}\right)^{d/2}
 \right].
 \label{eq:sharp-tail-second-moment}
\end{equation}
\end{corollary}

\begin{proof}
Chebyshev applied to \cref{thm:quenched-Abel-variance}, together with
\eqref{eq:Abel-mean} and the samplewise envelope
\eqref{eq:Abel-dominates-sharp}, proves
\eqref{eq:quenched-sharp-tail}.  The same envelope and
$\E\mathfrak A_L^2=(\E\mathfrak A_L)^2+
\Var\mathfrak A_L$ give \eqref{eq:sharp-tail-second-moment}.
\end{proof}

Define the effective number of diffusive soft modes
\begin{equation}
 m_{\mathrm{eff}}(L,\delta)
 =L^d\left(\frac{\delta}{j_+}\right)^{d/2}.
 \label{eq:effective-mode-count}
\end{equation}
The ratio of the fluctuation term in \eqref{eq:quenched-sharp-tail} to the
natural annealed scale $(\delta/j_+)^{d/2}$ is
$m_{\mathrm{eff}}^{-1/2}$.  Thus the fluctuation correction is lower order
for $\delta/j_+\gg L^{-2}$; whenever the mean has a nondegenerate
$\delta^{d/2}$ coefficient, this is relative self-averaging.  At
$\delta\asymp j_+L^{-2}$ only order-one effective modes are sampled and
order-one relative fluctuations are natural.  Below
\eqref{eq:finite-size-gap}, the sharp tail is exactly zero.


\begin{corollary}[Every-cutoff centered bound in the bounded-mode window]
\label{cor:bounded-effective-mode-sharp-variance}
Under the hypotheses of \cref{thm:quenched-Abel-variance}, define
\begin{equation}
 V_L(\delta)=j_+^2\abs A^4L^{-d}
 \left(\frac{\delta}{j_+}\right)^{d/2}.
 \label{eq:sharp-target-variance-scale-early}
\end{equation}
Then, at every deterministic cutoff,
\begin{equation}
 \boxed{
 \Var T_L(\delta)
 \le C\bigl(1+m_{\mathrm{eff}}(L,\delta)\bigr)V_L(\delta).}
 \label{eq:bounded-mode-centered-bound}
\end{equation}
Consequently, for each fixed $M<\infty$,
\begin{equation}
 m_{\mathrm{eff}}(L,\delta)\le M
 \quad\Longrightarrow\quad
 \Var T_L(\delta)\le C_MV_L(\delta).
 \label{eq:bounded-mode-centered-target}
\end{equation}
No cancellation around the mean is used here: the estimate is the elementary inequality $\Var T_L\le\E T_L^2$.  It has the centered target scale only when $m_{\mathrm{eff}}$ is bounded and becomes noninformative for the mesoscopic many-mode regime.  In particular, the desired centered nonlinear sharp-cutoff scale already holds
at every cutoff in any fixed acoustic window
$\delta\le C_0j_+L^{-2}$; below the deterministic gap the observable vanishes
exactly.
\end{corollary}

\begin{proof}
Since variance is bounded by the second moment, rewrite
\eqref{eq:sharp-tail-second-moment} using
$m_{\mathrm{eff}}=L^d(\delta/j_+)^{d/2}$:
\[
 j_+^2\abs A^4(\delta/j_+)^d
 =m_{\mathrm{eff}}(L,\delta)V_L(\delta).
\]
This proves \eqref{eq:bounded-mode-centered-bound} and its consequences.
\end{proof}

This observation removes the bottom of the diffusive window from the open
problem.  Genuine pointwise centering is needed only in the mesoscopic
many-mode regime $m_{\mathrm{eff}}\to\infty$.

\section{Exact centered theorem in the Born sector}

The sharp projector can be treated without smoothing in the frozen-spectrum
Born sector.  Put
\begin{equation}
 J_\eps=\Id+\eps\diag(\xi_e),
 \qquad \E\xi_e=0,
 \label{eq:Born-conductance}
\end{equation}
and let $L_0=BB^*$ be the clean torus Laplacian.  Since $Ba=0$, the exact
first-order drift is $g_\eps=\eps B\diag(\xi)a$.  Define
\begin{equation}
 T_L^{\mathrm B}(\delta)
 =\frac{\eps^2}{L^d}\xi^TM_\delta\xi,
 \qquad
 M_\delta=D_aB^*L_0^+\bm1_{(0,\delta]}(L_0)BD_a,
 \label{eq:Born-tail-quadratic-form}
\end{equation}
where $D_a=\diag(a_e)$.  If $\delta$ is not a positive eigenvalue of
$L_0$, then $T_L^{\mathrm B}(\delta)$ is the $\eps^2$ coefficient of the full
sharp tail as $\eps\to0$.  At a clean eigenvalue threshold the discontinuous
sharp functional need not possess a two-sided Taylor coefficient;
\eqref{eq:Born-tail-quadratic-form} is then the canonical frozen-spectrum Born
observable, and the theorem below concerns that observable exactly.

\begin{theorem}[Born sharp-tail concentration]
\label{thm:Born-sharp-concentration}
Let the $\xi_e$ be independent, centered, and sub-Gaussian with
$\norm{\xi_e}_{\psi_2}\le K$.  For $0<\delta\le1$,
\begin{align}
 \norm{M_\delta}&\le\abs A^2,
 &
 \norm{M_\delta}_{\mathrm F}^2
 &\le C_d\abs A^4L^d\delta^{d/2},
 \label{eq:Born-matrix-bounds}
\end{align}
and
\begin{equation}
 \boxed{
 \Prob\left[
 \abs{T_L^{\mathrm B}(\delta)-\E T_L^{\mathrm B}(\delta)}\ge r
 \right]
 \le2\exp\left[-c\min\left\{
 \frac{L^dr^2}{\eps^4K^4\abs A^4\delta^{d/2}},
 \frac{L^dr}{\eps^2K^2\abs A^2}
 \right\}\right].}
 \label{eq:Born-Hanson-Wright}
\end{equation}
In particular,
\begin{equation}
 \Var T_L^{\mathrm B}(\delta)
 \le C\eps^4K^4\abs A^4L^{-d}\delta^{d/2}.
 \label{eq:Born-variance}
\end{equation}
If $\sigma^2=\E\xi_e^2$, then the exact mean is
\begin{equation}
 \boxed{
 \E T_L^{\mathrm B}(\delta)
 =\frac{\eps^2\sigma^2}{L^d}
 \sum_{\substack{q\ne0\\\lambda(q)\le\delta}}
 \frac{\sum_{\mu=1}^dA_\mu^2\,4\sin^2(q_\mu/2)}{\sum_{\nu=1}^d4\sin^2(q_\nu/2)}.}
 \label{eq:Born-exact-mean}
\end{equation}
\end{theorem}

\begin{proof}
The operator inequality
$BD_a^2B^*\preceq\abs A^2BB^*$ gives the norm bound in
\eqref{eq:Born-matrix-bounds}.  The rank of $M_\delta$ is at most the number
of clean torus eigenvalues in $(0,\delta]$, which is bounded by
$C_dL^d\delta^{d/2}$ for $0<\delta\le1$.  Hence
$\norm{M_\delta}_{\mathrm F}^2\le
\operatorname{rank}(M_\delta)\norm{M_\delta}^2$.
The Hanson--Wright inequality \cite{RudelsonVershynin2013} applied to
\eqref{eq:Born-tail-quadratic-form} proves
\eqref{eq:Born-Hanson-Wright} and \eqref{eq:Born-variance}.
Finally, Fourier diagonalization of $L_0$ and
$\E\xi_e\xi_{e'}=\sigma^2\delta_{ee'}$ gives
\eqref{eq:Born-exact-mean}.
\end{proof}

\begin{remark}[Exact Born variance formula]
If $\kappa_4=\E\xi_e^4-3\sigma^4$, then
\begin{equation}
 \Var T_L^{\mathrm B}
 =\frac{\eps^4}{L^{2d}}
 \left[2\sigma^4\tr(M_\delta^2)
 +\kappa_4\sum_e(M_\delta)_{ee}^2\right].
 \label{eq:Born-exact-variance}
\end{equation}
This is the formula checked by the numerical audit.
\end{remark}

\section{Three closure routes: exact progress and sharp obstructions}
\label{sec:three-sharp-routes}

The smooth variance theorem leaves a precise question: can one transfer the
centered estimate to the discontinuous cumulative function
\eqref{eq:sharp-tail-definition}?  We now audit the three proposed routes.
Two of them yield exact new representations, and the first yields a
finite-volume Wegner theorem, but the analysis also shows why none of the
standard count-level inputs alone closes the fixed-cutoff theorem.

For a Borel set $I\subset(0,\infty)$, write
\begin{equation}
 \mu_{L,J}(I)
 =\frac1{L^d}\ip{g_L}{L_{J,L}^+\bm{1}_I(L_{J,L})g_L}.
 \label{eq:drift-Hminus1-measure}
\end{equation}
Thus $T_L(\delta)=\mu_{L,J}((0,\delta])$ and
$S_L(\delta,\eta)=\mu_{L,J}((\delta,(1+\eta)\delta])$.

\paragraph{Route I: a relative Wegner estimate.}
Assume in this paragraph that the edge variables are independent with common
density $\rho$ supported in $[j_-,j_+]$, and put
$\rho_\infty=\norm{\rho}_{\infty}$.  Let
\begin{equation}
 \mathcal N_L(I)=\operatorname{Tr}_{\bm{1}^\perp}\bm{1}_I(L_{J,L})
 \label{eq:positive-eigenvalue-count}
\end{equation}
be the positive-eigenvalue count.

\begin{theorem}[Finite discrete relative Wegner estimate]
\label{thm:relative-Wegner}
For every interval $I=(E,E+\ell]\subset(0,\infty)$,
\begin{equation}
 \boxed{
 \E\mathcal N_L(I)
 \le dL^d\rho_\infty j_+\frac{\ell}{E}.}
 \label{eq:relative-Wegner}
\end{equation}
Together with ellipticity and the clean-torus Weyl count,
\begin{equation}
 \E\mathcal N_L(I)
 \le C_dL^d\min\left\{
 \left(\frac{E+\ell}{j_-}\right)^{d/2},
 \rho_\infty j_+\frac{\ell}{E}
 \right\}.
 \label{eq:relative-Wegner-Weyl}
\end{equation}
\end{theorem}

\begin{proof}
Put $L_0=BB^*=\sum_eb_eb_e^*$ and choose the nondecreasing ramp
\[
 F(\lambda)=
 \begin{cases}
 0,&\lambda\le E,\\
 \lambda-E,&E<\lambda<E+\ell,\\
 \ell,&\lambda\ge E+\ell.
 \end{cases}
\]
Mollify the piecewise-linear ramp $F$ in arbitrarily small neighborhoods of its two corners, apply the preceding identity to the smooth monotone approximants, and pass to the limit by monotone convergence on both sides.  Since
$L_{J,L}\preceq j_+L_0$,
\[
 \mathcal N_L(I)
 \le E^{-1}\operatorname{Tr}[L_{J,L}\bm{1}_I(L_{J,L})]
 \le \frac{j_+}{E}\operatorname{Tr}[L_0\bm{1}_I(L_{J,L})].
\]
Functional differentiation gives
\[
 \sum_e\partial_{j_e}\operatorname{Tr} F(L_{J,L})
 =\operatorname{Tr}[F'(L_{J,L})L_0]
 =\operatorname{Tr}[\bm{1}_I(L_{J,L})L_0].
\]
Condition on every coefficient except $j_e=s$.  Since
$\partial_sL_{J^{(e,s)},L}=b_eb_e^*\succeq0$, every eigenvalue and hence
$s\mapsto\operatorname{Tr}F(L_{J^{(e,s)},L})$ is nondecreasing.  Together
with boundedness of the density this implies
\[
 \int\rho(s)\partial_s\operatorname{Tr} F(L_{J^{(e,s)},L})\dd s
 \le\rho_\infty
 \bigl[\operatorname{Tr} F(L_{j_+})-\operatorname{Tr} F(L_{j_-})\bigr].
\]
The endpoint matrices differ by the positive rank-one operator
$(j_+-j_-)b_eb_e^*$.  Rank-one interlacing and the fact that the range of $F$
has length $\ell$ give
\[
 0\le\operatorname{Tr} F(L_{j_+})-\operatorname{Tr} F(L_{j_-})\le\ell.
\]
Summing over the $dL^d$ positive-coordinate bonds proves
\eqref{eq:relative-Wegner}.  Finally,
$L_{J,L}\succeq j_-L_0$ and the lattice Weyl count give
\eqref{eq:relative-Wegner-Weyl}.
\end{proof}

The count theorem is the finite discrete analogue of divergence-form Wegner
estimates, for which eigenvalue lifting or gradient unique continuation is
needed in the continuum \cite{Dicke2021,DickeVeselic2023}.  Here the exact
bond covering $L_0=\sum_eb_eb_e^*$ and rank-one interlacing make the proof
self-contained.

It is nevertheless not a drift-weighted Wegner theorem.  The modal estimate is
only
\begin{equation}
 \frac1{L^d}\frac{\abs{\ip{g_L}{\psi_k}}^2}{\lambda_k}
 \le j_+\abs A^2,
 \label{eq:single-mode-coherent-bound}
\end{equation}
not $O(L^{-d})$.  A coherent soft mode can therefore carry an order-one
fraction of the whole $H^{-1}$ source mass.

\begin{proposition}[Normalized eigenvalue count does not control drift weight]
\label{prop:count-does-not-control-drift}
On the $N$-cycle, let
\begin{equation}
 j_{e,N}=1+\eps\cos(2\pi e/N),\qquad 0<\eps<1,
 \label{eq:modulated-cycle-second-use}
\end{equation}
and take $a_e=1$.  There are constants
$c_\eps,C_\eps>0$ such that, for all sufficiently large $N$,
\begin{equation}
 \boxed{
 \mu_N((0,C_\eps N^{-2}])\ge c_\eps,
 \qquad
 \frac1N\mathcal N_N((0,C_\eps N^{-2}])\longrightarrow0.}
 \label{eq:count-weight-obstruction}
\end{equation}
In particular, no realization-wise inequality
$\mu_{L,J}(I)\le C L^{-d}\mathcal N_L(I)$ can hold uniformly over bounded
elliptic conductances.
\end{proposition}

\begin{proof}
The total mass is the arithmetic-minus-harmonic conductance,
\[
 \mu_N((0,\infty))
 =1-\left(\frac1N\sum_ej_{e,N}^{-1}\right)^{-1}
 \longrightarrow1-\sqrt{1-\eps^2}>0.
\]
The first moment is explicit:
\[
 \int\lambda\dd\mu_N(\lambda)
 =\frac1N\norm{g_N}^2
 =2\eps^2\sin^2(\pi/N)
 \le2\pi^2\eps^2N^{-2}.
\]
Markov's inequality for the measure $\mu_N$ therefore places a fixed positive
fraction of its mass below $C_\eps N^{-2}$ once $C_\eps$ is large enough.
On the other hand,
$L_{J_N}\succeq(1-\eps)L_0$, and the clean cycle has only
$O_{\eps,C}(1)$ eigenvalues below $CN^{-2}$.  Dividing by $N$ proves
\eqref{eq:count-weight-obstruction}.
\end{proof}

This proposition prevents a shortcut through the ordinary IDS.  Recent
uniform concentration results for normalized eigenvalue counting functions
remain count-level statements \cite{KamperEtAl2026}; the needed theorem must
control current overlap with the low-energy subspace.

\paragraph{Route II: randomized cutoffs.}
Let $\Theta_\delta$ be independent of the conductances and exponentially
distributed with mean $\delta$.

\begin{theorem}[Abel tail as an exactly randomized sharp cutoff]
\label{thm:exponential-random-cutoff}
For every realization,
\begin{align}
 \boxed{\E_\Theta T_L(\Theta_\delta)=\mathfrak A_L(\delta),}
 \label{eq:Abel-randomized-mean}\\
 \boxed{
 \E_\Theta T_L(\Theta_\delta)^2
 =\iint\e^{-\max\{\lambda,\mu\}/\delta}
 \dd\mu_{L,J}(\lambda)\dd\mu_{L,J}(\mu)
 \le\mathfrak A_L(2\delta)^2.}
 \label{eq:Abel-randomized-second-moment}
\end{align}
Consequently, for $0<\delta\le j_+/2$,
\begin{equation}
 \E_{J,\Theta}T_L(\Theta_\delta)^2
 \le Cj_+^2\abs A^4\left[
 \left(\frac\delta{j_+}\right)^d
 +L^{-d}\left(\frac\delta{j_+}\right)^{d/2}
 \right].
 \label{eq:randomized-sharp-second-moment}
\end{equation}
The disorder variance of the conditional mean is exactly the Abel variance:
\begin{equation}
 \Var_J\bigl(\E_\Theta[T_L(\Theta_\delta)\mid J]\bigr)
 =\Var_J\mathfrak A_L(\delta),
 \label{eq:randomized-conditional-variance}
\end{equation}
and hence obeys \eqref{eq:quenched-Abel-variance}.
\end{theorem}

\begin{proof}
The survival probability is
$\Prob(\Theta_\delta\ge\lambda)=\e^{-\lambda/\delta}$, so Tonelli gives
\eqref{eq:Abel-randomized-mean}.  Applying the same identity twice gives the
double integral in \eqref{eq:Abel-randomized-second-moment}.  Since
$\max\{\lambda,\mu\}\ge(\lambda+\mu)/2$, its integrand is bounded by
$\e^{-\lambda/(2\delta)}\e^{-\mu/(2\delta)}$.  The final claims follow from
\eqref{eq:Abel-mean} and \eqref{eq:quenched-Abel-variance} at scale $2\delta$.
\end{proof}

Thus the Abel theorem is already an optimal centered disorder theorem for a
precisely defined randomized-frequency protocol.  It does not, however,
recover a deterministic threshold.  The total-variance identity
\begin{equation}
 \Var_{J,\Theta}T_L(\Theta_\delta)
 =\Var_J\mathfrak A_L(\delta)
 +\E_J\Var_\Theta(T_L(\Theta_\delta)\mid J)
 \label{eq:randomized-total-variance}
\end{equation}
contains threshold noise.  For the one-atom measure
$\mu=m\delta_{\{\delta\}}$, that term is
$m^2\e^{-1}(1-\e^{-1})$; if $m\asymp\delta^{d/2}$, it is of order
$\delta^d$, not $L^{-d}\delta^{d/2}$.  Broad randomization therefore cannot
by itself prove the fixed-cutoff CLT scale.  In fact this obstruction applies
to every positive heat-kernel mixture.

\begin{proposition}[No narrow cutoff from a positive heat mixture]
\label{prop:positive-heat-mixture-no-go}
Let a randomized-threshold survival function have the completely monotone form
\begin{equation}
 s(\lambda)=\int_0^\infty\e^{-t\lambda}\,\pi(\dd t),
 \qquad \pi((0,\infty))=1.
 \label{eq:positive-heat-mixture-survival}
\end{equation}
If the associated threshold $\Theta$ has finite second moment, then
\begin{equation}
 \boxed{\Var\Theta\ge(\E\Theta)^2.}
 \label{eq:positive-heat-mixture-CV}
\end{equation}
Thus no family of positive heat mixtures can concentrate its threshold in a
relative window $o(1)$ around its mean.
\end{proposition}

\begin{proof}
Conditionally on $t$, let $\Theta$ be exponential with rate $t$.
Then \eqref{eq:positive-heat-mixture-survival} is its survival function and
\[
 \E\Theta=\int t^{-1}\pi(\dd t),\qquad
 \E\Theta^2=2\int t^{-2}\pi(\dd t).
\]
Cauchy--Schwarz gives
$\int t^{-2}\dd\pi\ge(\int t^{-1}\dd\pi)^2$, proving
\eqref{eq:positive-heat-mixture-CV}.
\end{proof}

The proposition explains why positivity and narrowness conflict in a pure
heat-mixture strategy.  Erlang or gamma thresholds with large shape do
concentrate, but their survival functions contain polynomial times exponential
multipliers rather than positive mixtures of exponentials.  Exploiting them
would require uniform high-order semigroup-derivative estimates.

For a narrow randomization, let $U_{\delta,\eta}$ be uniform on
$[\delta,(1+\eta)\delta]$ and put
\begin{equation}
 \overline T_{L,\eta}(\delta)
 =\E_U[T_L(U_{\delta,\eta})\mid J].
 \label{eq:uniform-randomized-tail}
\end{equation}
Its spectral multiplier equals one below $\delta$, decreases linearly across
the shell, and vanishes above $(1+\eta)\delta$.  Hence, samplewise,
\begin{align}
 \boxed{0\le\overline T_{L,\eta}(\delta)-T_L(\delta)
 \le S_L(\delta,\eta),}
 \label{eq:uniform-randomized-sandwich}\\
 \boxed{
 \Var_U(T_L(U_{\delta,\eta})\mid J)
 \le\frac14S_L(\delta,\eta)^2.}
 \label{eq:uniform-randomized-variance}
\end{align}
Narrow randomization gives the desired smooth/sharp sandwich exactly, but it
also shows that the shell modulus cannot be avoided.

\paragraph{Route III: the integrated spectral variable.}
Define
\begin{equation}
 K_L(r)=\int_0^rT_L(s)\dd s
 =\int(r-\lambda)_+\dd\mu_{L,J}(\lambda).
 \label{eq:integrated-drift-spectral-measure}
\end{equation}
The multiplier is continuous at $\lambda=r$, and
\begin{equation}
 \boxed{
 \frac{K_L(r+h)-K_L(r)}h
 =\frac1h\int_r^{r+h}T_L(s)\dd s.}
 \label{eq:integrated-difference-quotient}
\end{equation}
For $h=\eta r$, the right-hand side is
$\overline T_{L,\eta}(r)$, so
\begin{equation}
 0\le\frac{K_L(r+h)-K_L(r)}h-T_L(r)
 \le S_L(r,h/r).
 \label{eq:integrated-shell-sandwich}
\end{equation}
Thus the integrated-variable and randomized-cutoff routes are the same
regularization in two languages.  Although pointwise shell control remains
open, its logarithmic-energy average closes exactly.

\begin{theorem}[Exact logarithmic shell identity and averaged shell bound]
\label{thm:logarithmic-shell-average}
For every finite positive measure $\mu$ on $(0,\infty)$, set
$S_\mu(r,\eta)=\mu((r,(1+\eta)r])$.  Then
\begin{equation}
 \boxed{
 \int_0^\infty S_\mu(r,\eta)^2\frac{\dd r}{r}
 =\iint
 \left[\log(1+\eta)-\abs{\log(\lambda/\lambda')}\right]_+
 \mu(\dd\lambda)\mu(\dd\lambda').}
 \label{eq:exact-log-shell-identity}
\end{equation}
Consequently, for the drift measure \eqref{eq:drift-Hminus1-measure},
$0<\eta\le1$, and $0<R\le j_+/(1+\eta)$,
\begin{equation}
 \boxed{
 \int_0^R\E S_L(r,\eta)^2\frac{\dd r}{r}
 \le C\eta j_+^2\abs A^4\left[
 \left(\frac R{j_+}\right)^d
 +L^{-d}\left(\frac R{j_+}\right)^{d/2}
 \right].}
 \label{eq:log-averaged-shell-bound}
\end{equation}
The same estimate holds for the integrated difference quotient:
\begin{equation}
 \int_0^R\E\left|\frac{K_L(r+\eta r)-K_L(r)}{\eta r}
 -T_L(r)\right|^2\frac{\dd r}{r}
 \le \textup{right-hand side of \eqref{eq:log-averaged-shell-bound}}.
 \label{eq:log-averaged-L2-differentiability}
\end{equation}
\end{theorem}

\begin{proof}
For fixed $\lambda$, the condition
$r<\lambda\le(1+\eta)r$ means that
$\log r$ lies in an interval of length $\log(1+\eta)$.  The intersection of
the two intervals associated with $\lambda$ and $\lambda'$ has length
$[\log(1+\eta)-|\log(\lambda/\lambda')|]_+$.  Tonelli proves
\eqref{eq:exact-log-shell-identity}.  Restricting $r\le R$ and discarding the
intersection geometry gives the deterministic bound
\begin{equation}
 \int_0^R S_L(r,\eta)^2\frac{\dd r}{r}
 \le\log(1+\eta)T_L((1+\eta)R)^2.
 \label{eq:localized-log-shell-deterministic}
\end{equation}
Take expectations, use \eqref{eq:sharp-tail-second-moment}, and absorb
$1+\eta\le2$ and $\log(1+\eta)\le\eta$.  Finally,
\eqref{eq:integrated-shell-sandwich} proves
\eqref{eq:log-averaged-L2-differentiability}.
\end{proof}


\begin{lemma}[Pointwise maximalization at the unavoidable floor]
\label{lem:log-shell-maximalization}
For every finite positive measure $\mu$, every $r>0$, and every $\eta>0$, put
$h=\log(1+\eta)$ and $\eta^{(2)}=(1+\eta)^2-1$.  Then
\begin{equation}
 \boxed{
 S_\mu(r,\eta)^2
 \le\frac1h\int_{r/(1+\eta)}^r
 S_\mu(s,\eta^{(2)})^2\frac{\dd s}{s}.}
 \label{eq:pointwise-log-shell-maximalization}
\end{equation}
Hence, for $0<\eta\le\sqrt2-1$ and
$r\le j_+/(1+\eta)^2$,
\begin{equation}
 \boxed{
 \E S_L(r,\eta)^2
 \le Cj_+^2\abs A^4\left[
 (r/j_+)^d+L^{-d}(r/j_+)^{d/2}\right].}
 \label{eq:pointwise-shell-floor}
\end{equation}
There is no exceptional cutoff in \eqref{eq:pointwise-shell-floor}, but the
bound has no vanishing power of $\eta$.
\end{lemma}

\begin{proof}
Push $\mu$ forward by $\lambda\mapsto x=\log\lambda$.  If
$y\in[x-h,x]$, then the interval $(x,x+h]$ is contained in
$(y,y+2h]$.  Thus its mass squared is bounded by the average of the wider-shell
mass squared over $y\in[x-h,x]$, which is
\eqref{eq:pointwise-log-shell-maximalization}.  Take expectations and apply
\eqref{eq:log-averaged-shell-bound} with relative width
$\eta^{(2)}\le1$ and upper endpoint $r$.  Since
$\eta^{(2)}/\log(1+\eta)$ is uniformly bounded on
$(0,\sqrt2-1]$, this gives \eqref{eq:pointwise-shell-floor}.
\end{proof}

The lemma is a direct conclusion obtainable from positivity and the
current second-moment input.  In the bounded-mode window its right-hand
side is the target variance scale by
\cref{cor:bounded-effective-mode-sharp-variance}; in the mesoscopic window the
first term is larger by $m_{\mathrm{eff}}$.  The acoustic crossing model below
shows why a volume-uniform improvement by a factor $\eta^\alpha$ cannot be
expected at arbitrarily thin prescribed shells.

In particular, if $R_*$ is independent and log-uniform on $[R/2,R]$, then
\begin{equation}
 \E_{J,R_*}S_L(R_*,\eta)^2
 \le C\eta j_+^2\abs A^4\left[
 (R/j_+)^d+L^{-d}(R/j_+)^{d/2}
 \right].
 \label{eq:log-uniform-random-center-shell}
\end{equation}
Thus the desired shell modulus is already true for a typical cutoff in every
dyadic energy window.  What is missing is a bound uniform at a prescribed
cutoff.


\begin{corollary}[Exceptional set of spectral cutoffs]
\label{cor:exceptional-cutoff-set}
Fix $0<\beta<1$, $0<\eta\le1$, and put
\begin{equation}
 Q_L(R)=j_+^2\abs A^4\left[
 (R/j_+)^d+L^{-d}(R/j_+)^{d/2}
 \right].
 \label{eq:dyadic-shell-scale}
\end{equation}
There is a constant $C$ such that the bad set
\begin{equation}
 \mathcal B_{L,R}(\eta,\beta)
 =\left\{r\in[R/2,R]:
 \E S_L(r,\eta)^2>C\eta^\beta Q_L(R)\right\}
 \label{eq:bad-cutoff-set}
\end{equation}
obeys
\begin{equation}
 \boxed{
 \int_{\mathcal B_{L,R}(\eta,\beta)}\frac{\dd r}{r}
 \le C\eta^{1-\beta}.}
 \label{eq:bad-cutoff-log-measure}
\end{equation}
Consequently, let $(L_k)$ be any deterministic volume sequence and let
$\eta_k\downarrow0$ satisfy
$\sum_k\eta_k^{1-\beta}<\infty$.  Then for logarithmically almost every
$r\in[R/2,R]$,
\begin{equation}
 \E S_{L_k}(r,\eta_k)^2
 \le C\eta_k^\beta Q_{L_k}(R)
 \label{eq:almost-every-cutoff-shell}
\end{equation}
for all sufficiently large $k$.  The same conclusion holds for the
$L^2$ error between $T_{L_k}(r)$ and the collapsing integrated difference
quotient with $h_k=\eta_kr$.
\end{corollary}

\begin{proof}
Restrict \eqref{eq:log-averaged-shell-bound} to $[R/2,R]$ and apply Markov's
inequality with respect to logarithmic measure.  This gives
\eqref{eq:bad-cutoff-log-measure}.  The second statement follows from the
Borel--Cantelli lemma on the finite measure space
$([R/2,R],\dd r/r)$; no independence in $k$ is needed.  The final claim uses
\eqref{eq:integrated-shell-sandwich}.
\end{proof}

Thus the pointwise shell problem is already solved outside a quantitatively
small deterministic set of cutoff energies, and along arbitrary volume
sequences it is solved for almost every energy.  The remaining theorem asks
for a bound at a prescribed cutoff with no exceptional set.

Order the $dL^d$ conductances, let $(\mathcal F_m)$ be their reveal filtration,
and put
\begin{equation}
 Y_{r,h}=\frac{K_L(r+h)-K_L(r)}h,
 \qquad
 D_m(r,h)=\E[Y_{r,h}\mid\mathcal F_m]
 -\E[Y_{r,h}\mid\mathcal F_{m-1}].
 \label{eq:integrated-Doob-increments}
\end{equation}
The exact Doob identity is
\begin{equation}
 \Var Y_{r,h}=\sum_{m=1}^{dL^d}\E\abs{D_m(r,h)}^2.
 \label{eq:integrated-Doob-identity}
\end{equation}


\section{Pointwise rank-one geometry, participation coarea, and the correct martingale target}
\label{sec:pointwise-rank-one-coarea}

The logarithmic shell theorem leaves open whether its exceptional set can be
removed by a pointwise maximal estimate.  We now resolve the finite-dimensional
geometry of that question.  The first conclusion is favorable: varying one
bond cannot produce multiple crossings of a fixed cutoff.  The second is a
warning: the resulting one-bond Efron--Stein surface term is generally stronger
than the centered theorem and can lose a square-root volume factor at an
individual acoustic eigenvalue.  This identifies direct Doob smoothing in
\emph{disorder space}, rather than a volume-uniform spectral shell modulus, as
the correct remaining route.

Write $N=L^d$, $b_e=B\mathbf e_e$, and, after deleting the $e$th bond,
\begin{equation}
 L_e^{(0)}=\sum_{f\ne e}j_fb_fb_f^*,
 \qquad
 g_e^{(0)}=\sum_{f\ne e}j_fa_fb_f.
 \label{eq:leave-one-edge-operator-source}
\end{equation}
For $t\in[j_-,j_+]$ put
\begin{equation}
 L_e(t)=L_e^{(0)}+tb_eb_e^*,
 \qquad
 g_e(t)=g_e^{(0)}+ta_eb_e.
 \label{eq:one-edge-rank-one-path}
\end{equation}
All operators below are restricted to $\mathbf1^\perp$.

\begin{lemma}[Single-crossing lemma]
\label{lem:single-crossing}
For $t_1<t_2$ and every $r>0$,
\begin{equation}
 0\le
 \tr\bm1_{(0,r]}(L_e(t_1))-
 \tr\bm1_{(0,r]}(L_e(t_2))
 \le1.
 \label{eq:single-crossing-count}
\end{equation}
Thus a complete one-bond path crosses a prescribed positive energy at most
once, counting multiplicity.
\end{lemma}

\begin{proof}
The perturbation
$L_e(t_2)-L_e(t_1)=(t_2-t_1)b_eb_e^*$ is positive and rank one.  Monotonicity
and rank-one Cauchy interlacing imply that the counting function below $r$ can
only decrease and can decrease by at most one.
\end{proof}

The unique crossing has explicit leave-one-edge coordinates.  For
$r\notin\sigma(L_e^{(0)}|_{\mathbf1^\perp})$, define
\begin{align}
 m_e(r)&=\ip{b_e}{(L_e^{(0)}-r)^{-1}b_e},
 &
 \dot m_e(r)&=\ip{b_e}{(L_e^{(0)}-r)^{-2}b_e},
 \label{eq:Krein-m-functions}\\
 h_e(r)&=\ip{b_e}{(L_e^{(0)}-r)^{-1}g_e^{(0)}}.
 \label{eq:Krein-h-function}
\end{align}

\begin{proposition}[Krein coordinates for the unique crossing]
\label{prop:Krein-crossing-coordinates}
A crossing occurs in the ellipticity interval precisely when
\begin{equation}
 t_e(r)=-\frac1{m_e(r)}\in(j_-,j_+).
 \label{eq:Krein-crossing-location}
\end{equation}
Because the admissible conductance is positive, such a crossing necessarily has $m_e(r)<0$; equivalently, $r$ lies above at least one eigenvalue of the deleted-bond operator on $\mathbf1^\perp$.
At that crossing, a normalized eigenvector, the eigenvalue speed, and the jump
of the sharp tail are
\begin{align}
 \psi_e(r)
 &=\frac{(L_e^{(0)}-r)^{-1}b_e}{\sqrt{\dot m_e(r)}},
 \label{eq:Krein-crossing-vector}\\
 \left.\partial_t\lambda_e(t)\right|_{t=t_e(r)}
 &=\frac{m_e(r)^2}{\dot m_e(r)},
 \label{eq:Krein-crossing-speed}\\
 q_e(r)
 &=\frac1{Nr}\abs{\ip{g_e(t_e(r))}{\psi_e(r)}}^2
 =\frac{\abs{h_e(r)-a_e}^2}{Nr\dot m_e(r)}.
 \label{eq:Krein-crossing-jump}
\end{align}
As $t$ increases through $t_e(r)$, the function
$T_L(r;J_{\setminus e},t)$ has a single downward jump of size $q_e(r)$.
Moreover, for every $z\notin\sigma(L_e(t))$,
\begin{align}
 &\ip{g_e(t)}{(L_e(t)-z)^{-1}g_e(t)}
 \nonumber\\
 &\quad=
 \ip{g_e^{(0)}}{(L_e^{(0)}-z)^{-1}g_e^{(0)}}
 +\frac{t\bigl(a_e^2m_e(z)t+2a_eh_e(z)-h_e(z)^2\bigr)}
 {1+tm_e(z)}.
 \label{eq:affine-source-Krein-resolvent}
\end{align}
\end{proposition}

\begin{proof}
The matrix determinant lemma gives
$\det(L_e(t)-r)=\det(L_e^{(0)}-r)(1+tm_e(r))$, which proves
\eqref{eq:Krein-crossing-location}.  Applying $L_e(t)-r$ to
$(L_e^{(0)}-r)^{-1}b_e$ gives $(1+tm_e(r))b_e$, and normalization yields
\eqref{eq:Krein-crossing-vector}.  Hellmann--Feynman gives
$\partial_t\lambda=|\langle b_e,\psi_e\rangle|^2=m_e^2/\dot m_e$.
At $t=-1/m_e$,
\[
 \ip{g_e(t)}{\psi_e}
 =\frac{h_e+ta_em_e}{\sqrt{\dot m_e}}
 =\frac{h_e-a_e}{\sqrt{\dot m_e}},
\]
which proves \eqref{eq:Krein-crossing-jump}.  Finally, the
Sherman--Morrison formula, together with the affine dependence of $g_e(t)$,
gives \eqref{eq:affine-source-Krein-resolvent} by direct expansion.
\end{proof}

Let $(\lambda_k,\psi_k)$ denote the positive eigenpairs of $L_{J,L}$ and set
\begin{equation}
 q_k(J)=\frac1{N\lambda_k}\abs{\ip{g_L}{\psi_k}}^2,
 \qquad
 \Pi_L(I)=\E\sum_kq_k(J)^2\bm1_I(\lambda_k).
 \label{eq:participation-measure}
\end{equation}
The measure $\Pi_L$ is the diagonal, or inverse-participation, part of the
second moment of the drift spectral measure.  It is exactly the coefficient
of the linear small-shell term; off-diagonal pairs contribute only at higher
order when the disorder-averaged spectral law is regular.

\begin{theorem}[Participation coarea formula]
\label{thm:participation-coarea}
Assume that the edge law has density $\rho$ on $[j_-,j_+]$ and that the
positive spectrum is simple almost surely.  Then $\Pi_L$ is absolutely
continuous on $(0,\infty)$.  For almost every $r>0$, its density
$\pi_L(r)=\dd\Pi_L/\dd r$ satisfies
\begin{equation}
 \boxed{
 r\pi_L(r)=
 \sum_e\E_{\setminus e}\left[
 t_e(r)\rho(t_e(r))q_e(r)^2
 \bm1_{\{t_e(r)\in(j_-,j_+)\}}
 \right].}
 \label{eq:participation-coarea-density}
\end{equation}
Equivalently, the explicit leave-one-edge Green-function target is
\begin{equation}
 r\pi_L(r)=\frac1{N^2r^2}
 \sum_e\E_{\setminus e}\left[
 t_e(r)\rho(t_e(r))
 \frac{\abs{h_e(r)-a_e}^4}{\dot m_e(r)^2}
 \bm1_{\{t_e(r)\in(j_-,j_+)\}}
 \right].
 \label{eq:participation-Green-target}
\end{equation}
In particular,
\begin{equation}
 \E\sum_{r<\lambda_k\le(1+\eta)r}q_k^2
 =\int_r^{(1+\eta)r}\pi_L(s)\dd s.
 \label{eq:diagonal-shell-participation}
\end{equation}
\end{theorem}

\begin{proof}
Every positive eigenvalue is homogeneous of degree one in the conductances, so
Euler's identity and Hellmann--Feynman give
\begin{equation}
 \sum_e\frac{j_e\partial_{j_e}\lambda_k}{\lambda_k}=1.
 \label{eq:eigenvalue-Euler-identity}
\end{equation}
For a compactly supported test function $f$,
\[
 \int f(r)\Pi_L(\dd r)
 =\sum_e\E\sum_k
 \frac{j_e\partial_{j_e}\lambda_k}{\lambda_k}
 q_k^2f(\lambda_k).
\]
Condition on all bonds except $j_e=t$.  The eigenvalue branches are monotone,
and \cref{lem:single-crossing} says that at a fixed $r$ there is at most one
crossing over the full $t$ interval.  Changing variables from $t$ to the
crossing eigenvalue cancels $\partial_t\lambda_k$ and gives
\[
 \int f(r)\frac1r\E_{\setminus e}
 \left[t_e(r)\rho(t_e(r))q_e(r)^2
 \bm1_{\{t_e(r)\in(j_-,j_+)\}}\right]\dd r.
\]
Summing over $e$ proves \eqref{eq:participation-coarea-density};
\eqref{eq:participation-Green-target} follows from
\eqref{eq:Krein-crossing-jump}.
\end{proof}

The coarea formula also gives the exact one-bond surface term in the product
variance inequality.  Let $F_\beta(t)=\beta([j_-,t])$, and, conditionally on
$J_{\setminus e}$, put
\begin{equation}
 H_e(t)=T_L(r;J_{\setminus e},t),
 \qquad
 \widetilde H_e(t)=H_e(t)+q_e(r)\bm1_{\{t\ge t_e(r)\}},
 \label{eq:regularized-one-edge-sharp-path}
\end{equation}
with $q_e=0$ when there is no crossing.  The function $\widetilde H_e$ is
continuous and piecewise analytic.

\begin{theorem}[Exact Efron--Stein surface decomposition]
\label{thm:Efron-Stein-surface-decomposition}
Suppose the one-edge law satisfies the derivative Poincar\'e inequality
\eqref{eq:one-edge-Poincare}.  Define
\begin{align}
 \mathcal R_L(r)
 &=\sum_e\E_{\setminus e}\int
 \abs{\partial_t\widetilde H_e(t)}^2\,\beta(\dd t),
 \label{eq:regular-one-edge-sensitivity}\\
 \mathcal J_L(r)
 &=\sum_e\E_{\setminus e}\left[
 q_e(r)^2F_\beta(t_e(r))(1-F_\beta(t_e(r)))
 \bm1_{\{t_e(r)\in(j_-,j_+)\}}
 \right].
 \label{eq:jump-one-edge-sensitivity}
\end{align}
Then
\begin{equation}
 \boxed{
 \Var T_L(r)\le2C_\beta\mathcal R_L(r)+2\mathcal J_L(r).}
 \label{eq:Efron-Stein-surface-bound}
\end{equation}
If
\begin{equation}
 \kappa_\beta=
 \operatorname*{ess\,sup}_{j_-<t<j_+}
 \frac{F_\beta(t)(1-F_\beta(t))}{t\rho(t)}<\infty,
 \label{eq:one-edge-isoperimetric-ratio}
\end{equation}
then
\begin{equation}
 \boxed{\mathcal J_L(r)\le\kappa_\beta\,r\pi_L(r).}
 \label{eq:jump-participation-density-bound}
\end{equation}
For the uniform law on $[j_-,j_+]$,
$\kappa_\beta\le(j_+-j_-)/(4j_-)$.
\end{theorem}

\begin{proof}
Tensorization gives
$\Var T_L(r)\le\sum_e\E_{\setminus e}\Var_\beta H_e$.
Since
$H_e=\widetilde H_e-q_e\bm1_{\{t\ge t_e\}}$,
\[
 \Var_\beta H_e
 \le2\Var_\beta\widetilde H_e
 +2q_e^2F_\beta(t_e)(1-F_\beta(t_e)).
\]
Apply the one-dimensional Poincar\'e inequality to $\widetilde H_e$ and sum.
Equation \eqref{eq:jump-participation-density-bound} follows term by term from
\eqref{eq:participation-coarea-density} and
\eqref{eq:one-edge-isoperimetric-ratio}.
\end{proof}

This theorem removes the previously stated ``multiple crossing'' obstruction:
there is only one crossing channel.  It also shows that the naive pointwise
target
\begin{equation}
 r\pi_L(r)\stackrel{?}{\lesssim}
 j_+^2\abs A^4N^{-1}(r/j_+)^{d/2}
 \label{eq:too-strong-pointwise-participation-target}
\end{equation}
is an Efron--Stein sufficient condition, not a necessary consequence of the
centered variance theorem.  In fact, at a cutoff passing through an individual
acoustic eigenvalue, \eqref{eq:too-strong-pointwise-participation-target} is
not the scale predicted by eigenvalue central-limit fluctuations.

\begin{proposition}[Acoustic-CLT obstruction to a uniform maximal shell modulus]
\label{prop:acoustic-CLT-maximal-obstruction}
Let $N\to\infty$, let $r_N>0$, and consider a single positive modal
contribution
\begin{equation}
 \lambda_N=r_N\exp(Z_N/\sqrt N),
 \qquad q_N=Q_N/N.
 \label{eq:acoustic-one-mode-scaling}
\end{equation}
Assume $(Z_N,Q_N)\Rightarrow(Z,Q)$, the family $Q_N^2$ is uniformly
integrable, and for some $a>0$,
\begin{equation}
 \E\left[Q^2\bm1_{\{0<Z\le a\}}\right]>0.
 \label{eq:nondegenerate-acoustic-crossing}
\end{equation}
Put $\eta_N=\exp(a/\sqrt N)-1$.  Then the one-mode shell obeys
\begin{equation}
 N^2\E\left[q_N^2
 \bm1_{\{r_N<\lambda_N\le(1+\eta_N)r_N\}}\right]
 \longrightarrow
 \E\left[Q^2\bm1_{\{0<Z\le a\}}\right]>0.
 \label{eq:acoustic-shell-nonvanishing}
\end{equation}
Consequently, for every $\alpha>0$, no volume-uniform estimate of the form
\begin{equation}
 \E S_N(r_N,\eta)^2\le C\eta^\alpha N^{-2}
 \label{eq:false-uniform-Holder-shell}
\end{equation}
can hold at all $\eta$, even though the centered one-mode tail has variance
of order $N^{-2}$.  In the explicit Gaussian case
$Z\sim\mathcal N(0,1)$ and $Q\equiv q_*$,
\begin{equation}
 r_N\pi_N(r_N)=\frac{q_*^2}{\sqrt{2\pi}}N^{-3/2},
 \qquad
 \Var\left[\frac{q_*}{N}\bm1_{\{Z\le0\}}\right]
 =\frac{q_*^2}{4}N^{-2}.
 \label{eq:Gaussian-acoustic-density-versus-variance}
\end{equation}
Thus the pointwise one-bond surface density can be larger than the centered
variance scale by $\sqrt N$.
\end{proposition}

\begin{proof}
The shell event is exactly $\{0<Z_N\le a\}$.  Uniform integrability and weak
convergence give \eqref{eq:acoustic-shell-nonvanishing}.  Since
$\eta_N^\alpha\sim a^\alpha N^{-\alpha/2}$,
\eqref{eq:false-uniform-Holder-shell} would force the positive limit in
\eqref{eq:acoustic-shell-nonvanishing} to vanish.  The Gaussian formulas
follow from the density of $Z/\sqrt N$ at the origin and from the variance of
a fair Bernoulli variable.
\end{proof}

The proposition is an exact surrogate statement, not an assertion that a
specific torus eigenpair already satisfies the required joint limit.  It is,
however, the scaling expected from quantitative eigenvalue fluctuation theory
for isolated homogenized elliptic modes \cite{Duerinckx2022Eigenvalue}; large-scale acoustic dispersive estimates give complementary context \cite{DuerinckxGloria2025}.  The
finite-torus audit in the supplementary package finds
$\operatorname{sd}(\log\lambda)\simeq N^{-1/2}$,
$r\pi_L(r)\simeq N^{-3/2}$, and
$\Var T_L(r)\simeq N^{-2}$ across the first acoustic multiplet in two
dimensions.  These diagnostics do not enter the proof.

The correct closure statement must therefore allow a volume-dependent
spectral smoothing scale, or avoid spectral smoothing altogether.

\begin{theorem}[Diagonal integrated-martingale criterion]
\label{thm:diagonal-integrated-martingale-criterion}
Let $r_L$ lie in the diffusive window and define
\begin{equation}
 V_L(r_L)=j_+^2\abs A^4N^{-1}(r_L/j_+)^{d/2}.
 \label{eq:target-centered-variance-scale}
\end{equation}
Suppose there is a sequence $h_L\downarrow0$ such that
\begin{align}
 \E S_L(r_L,h_L/r_L)^2&=o(V_L(r_L)),
 \label{eq:diagonal-shell-error}\\
 \sum_m\E\abs{D_m(r_L,h_L)}^2&\le C V_L(r_L).
 \label{eq:diagonal-integrated-QV}
\end{align}
Then
\begin{equation}
 \Var T_L(r_L)\le(C+o(1))V_L(r_L).
 \label{eq:diagonal-centered-conclusion}
\end{equation}
No volume-uniform limit $h\downarrow0$ is required.
\end{theorem}

\begin{proof}
By \eqref{eq:integrated-shell-sandwich},
$\|Y_{r_L,h_L}-T_L(r_L)\|_2=o(V_L^{1/2})$.
The exact Doob identity \eqref{eq:integrated-Doob-identity} and the triangle
inequality for centered $L^2$ norms then give
\eqref{eq:diagonal-centered-conclusion}.
\end{proof}

There is an even cleaner fixed-cutoff formulation.  Define the sharp Doob
increments
\begin{equation}
 D_m^\sharp(r)=\E[T_L(r)\mid\mathcal F_m]
 -\E[T_L(r)\mid\mathcal F_{m-1}].
 \label{eq:sharp-Doob-increments}
\end{equation}
Then
\begin{equation}
 \boxed{
 \Var T_L(r)=\sum_m\E\abs{D_m^\sharp(r)}^2.}
 \label{eq:sharp-Doob-identity}
\end{equation}
The full centered theorem is therefore equivalent to a direct pointwise bound
on this quadratic variation.  Unlike Efron--Stein, the Doob increments average
each crossing over all unrevealed conductances and can capture the acoustic
central-limit smoothing responsible for the missing $N^{-1/2}$ factor.

A useful exact representation of that smoothing is provided by the reversible
noise semigroup on disorder space.  Let
\begin{equation}
 \mathscr G=\sum_e(\E_e-\Id),
 \qquad
 P_t=\e^{t\mathscr G},
 \label{eq:product-resampling-semigroup}
\end{equation}
where $\E_e$ averages only the $e$th conductance.  The product law is
reversible and
\begin{equation}
 \boxed{
 \Var F=2\int_0^\infty
 \sum_e\E\Var_e(P_tF)\dd t.}
 \label{eq:disorder-noise-variance-representation}
\end{equation}
Thus it is sufficient to prove, at the fixed physical cutoff,
\begin{equation}
 \sum_e\E\Var_e(P_tT_L(r))
 \le C V_L(r)\vartheta(t),
 \qquad
 \int_0^\infty\vartheta(t)\dd t<\infty.
 \label{eq:noise-energy-target}
\end{equation}
This route never differentiates the spectral projector in $r$ and has no
exceptional cutoff set.

The expected short-time singularity in \eqref{eq:noise-energy-target} is
integrable and is already visible in the exact Gaussian half-space model.  If
$X\in\R^M$ is standard Gaussian,
$Z=M^{-1/2}\sum_iX_i$, and $F=q\bm1_{\{Z\le0\}}$, then the
Ornstein--Uhlenbeck semigroup satisfies
\begin{equation}
 P_tF=q\,\Phi\left(
 -\frac{\e^{-t}Z}{\sqrt{1-\e^{-2t}}}\right).
 \label{eq:Gaussian-halfspace-semigroup}
\end{equation}
Writing $a_t=\e^{-t}/\sqrt{1-\e^{-2t}}$ gives
\begin{equation}
 \E\abs{\nabla P_tF}^2
 =\frac{q^2a_t^2}{2\pi\sqrt{1+2a_t^2}}
 \sim\frac{q^2}{4\pi\sqrt t}
 \quad(t\downarrow0),
 \label{eq:Gaussian-halfspace-noise-energy}
\end{equation}
while it decays exponentially for large $t$, and exactly
\begin{equation}
 2\int_0^\infty\E\abs{\nabla P_tF}^2\dd t
 =\frac{q^2}{4}=\Var F.
 \label{eq:Gaussian-halfspace-variance-closure}
\end{equation}
The $t^{-1/2}$ boundary layer is the disorder-space replacement for the false
volume-uniform shell modulus.  The new sharp target is to prove the analogue
of \eqref{eq:Gaussian-halfspace-noise-energy} for the low-energy
current-weighted random-conductance functional, using the coarea coordinates
\eqref{eq:participation-Green-target} together with conditional homogenization
of the unrevealed environment.


\begin{problem}[Pointwise disorder-noise energy]
\label{prob:pointwise-disorder-noise-energy}
Prove the fixed-cutoff estimate
\begin{equation}
 \int_0^\infty\sum_e\E\Var_e(P_tT_L(r))\dd t
 \le Cj_+^2\abs A^4N^{-1}(r/j_+)^{d/2}
 \label{eq:integrated-noise-energy-goal}
\end{equation}
uniformly for $cL^{-2}\le r/j_+\le\delta_0$.  A stronger, scale-resolved
version would establish \eqref{eq:noise-energy-target} with
$\vartheta(t)\lesssim t^{-1/2}$ for $t\le1$ and exponential decay for
$t\ge1$.  The coarea theorem identifies the exact crossing surface, while the
acoustic obstruction shows why the unrevealed-environment smoothing in $P_t$
or in the direct Doob martingale is essential.  Proving
\eqref{eq:integrated-noise-energy-goal} closes the centered nonlinear theorem
without any exceptional set of physical cutoff energies.
\end{problem}





\section{Numerical audit}
\label{sec:numerics}

The Paper~II archive retains seven independent programs.  They verify the Abel heat and parabolic-energy identities, the exact Duhamel bond derivative, the Born Fourier formulas and Hanson--Wright normalization, the relative Wegner and randomized-cutoff identities, the logarithmic shell formula, the single-crossing/Krein/coarea relations, and the Gaussian and Bernoulli disorder-noise boundary layers.  The exact Abel Duhamel/spectral derivative discrepancy is $2.17\times10^{-19}$, the finite-difference sensitivity error is $1.05\times10^{-12}$, the one-crossing eigenvalue error is $6.66\times10^{-16}$, and the exact Gaussian noise-energy integral closes to $7.92\times10^{-14}$.  Monte Carlo coarea and low-mode scaling tables are explicitly labeled diagnostics rather than proof inputs.

\section{Theorem boundary}
\label{sec:boundary}

The positive Abel envelope has the predicted effective-mode variance under the stated one-bond Poincar\'e and period-uniform GNO hypotheses.  The physical sharp tail inherits a one-sided high-probability bound.  At bounded effective-mode count, the target variance scale follows from the uncentered second moment; genuine centering remains necessary in the mesoscopic window.  The Born sharp observable has a complete centered concentration theorem.

The three-route audit closes several tempting shortcuts.  Counts alone do not control drift weight; positive heat mixtures cannot create a narrow threshold; logarithmic shell averaging does not imply a vanishing pointwise shell modulus; and a single bond creates at most one cutoff crossing.  The rank-one and coarea formulas isolate the exact surface term that direct Efron--Stein overcounts near an acoustic mode.

\begin{problem}[Fixed-cutoff disorder-noise theorem]
\label{prob:fixed-cutoff-noise}
Prove, uniformly for $cL^{-2}\le r/j_+\le\delta_0$ in the mesoscopic regime $m_{\rm eff}(L,r)\to\infty$,
\[
 2\int_0^\infty\sum_e\E\Var_e(P_tT_L(r))\dd t
 \le Cj_+^2\abs A^4L^{-d}(r/j_+)^{d/2}.
\]
A scale-resolved estimate with an integrable $t^{-1/2}$ short-time profile would close the centered nonlinear theorem at every prescribed cutoff without an exceptional energy set and without a false uniform maximal-shell hypothesis.
\end{problem}

\section{Conclusion}

The smooth and weak-disorder sectors now have complete fluctuation theorems at the effective diffusive-mode scale.  For the nonlinear sharp cutoff, the obstruction is not an unidentified eigenvalue crossing: rank-one geometry makes the crossing unique and explicit.  Nor is it a small exceptional set of physical frequencies: pointwise maximalization reaches the unavoidable uncentered floor at every cutoff.  The missing mechanism is averaging in disorder space.  The exact product-noise representation identifies the final theorem as an integrable fixed-cutoff quadratic-variation estimate, with the acoustic half-space model predicting the correct $t^{-1/2}$ boundary layer.


\appendix

\section{Periodic gradient propagator: definitions, imported estimates, and proof}
\label{app:gradient-propagator}

Normalize $j_+=1$ and put $\lambda_*=j_-/j_+$.  For a fixed periodic environment $J$, define the torus heat kernel by
\begin{equation}
 G_L(t,J,x,y)=\bigl(\e^{-tL_{J,L}}\delta_y\bigr)(x),
 \qquad x,y\in\Lambda_L.
 \label{eq:torus-Green-definition}
\end{equation}
It solves $\partial_tG_L+L_{J,L}G_L=0$ with initial datum $\delta_y$ and conserves total mass.  Gradients annihilate its constant long-time component, so no mean-zero subtraction is required in the estimates below.  Let
\begin{equation}
 \omega_{L,t}(y)=\left(1+\frac{\dist(y,L\Z^d)^2}{1+t}\right)^{1/2},
 \label{eq:GNO-weight-definition}
\end{equation}
where $\dist(y,L\Z^d)$ is the Euclidean distance of a representative of $y$ to the nearest period point.

We use two inputs from the long version of Gloria--Neukamm--Otto \cite{GloriaNeukammOtto2013Long}; the corresponding quantitative homogenization framework is published in \cite{GloriaNeukammOtto2015}.  The period-uniform statements below are taken from the numbered results of the long version.

For completeness, we spell out the periodic vertical hypothesis.  Let $Q_L=\{0,\ldots,L-1\}^d$, regard the $d$ outgoing coefficients at a site as one coefficient block, and set
\begin{equation}
 \partial_{L,y}F
 =F-\E\bigl[F\mid\{j(x):x\in Q_L,\ x\ne y\}\bigr].
 \label{eq:periodic-vertical-derivative}
\end{equation}
The condition $\mathrm{SG}_L(\rho)$ is
\begin{equation}
 \Var F\le\rho^{-1}\sum_{y\in Q_L}\E\abs{\partial_{L,y}F}^2
 \label{eq:periodic-SG-definition}
\end{equation}
for every square-integrable mean-zero $F$.  The product law of independent site blocks satisfies this inequality with a constant independent of $L$ by tensorization.  For a vector-valued random field $\xi$, write
\begin{equation}
 \norm{\xi}_{\ell_y^1L^{2p}}
 :=\sum_{y\in Q_L}\norm{\partial_{L,y}\xi}_{L^{2p}(\Omega_L)}.
 \label{eq:vertical-source-norm}
\end{equation}

\begin{theorem}[Imported divergence-semigroup estimate]
\label{thm:GNO-divergence-import}
Assume the stationary $L$-periodic ensemble satisfies $\mathrm{SG}_L(\rho)$ with $\rho$ independent of $L$.  Theorem~2 under alternative~(69b) supplies $p_0=p_0(d,\lambda_*)<\infty$ such that, for $p\ge p_0$ and $u(t)=\e^{-t\mathscr L_L}D^*\xi$,
\begin{equation}
 \norm{u(t)}_{L^{2p}(\Omega_L)}
 \le C_p(1+t)^{-d/4-1/2}\norm{\xi}_{\ell_y^1L^{2p}},
 \label{eq:GNO-Theorem2-specialized}
\end{equation}
where $C_p=C(p,d,\lambda_*,\rho)$ has no dependence on $L$.  Lower finite moments follow by Jensen.  For the local drift $d_A=D^*(j(0)A)$, Remark~11 and equation~(70) give
\begin{equation}
 \norm{\e^{-t\mathscr L_L}d_A}_{L^p(\Omega_L)}
 \le C_p\abs A(1+t)^{-d/4-1/2}.
 \label{eq:GNO-drift-moment}
\end{equation}
\end{theorem}

\begin{theorem}[Imported periodic Green-gradient estimate]
\label{thm:GNO-Green-import}
Theorem~3(b), with the weight in its equation~(74), supplies an exponent $q>1$ depending only on $d,\lambda_*$ such that for every $\alpha<\infty$,
\begin{equation}
 \left(\sum_{y\in\Lambda_L}
 [\omega_{L,t}(y)^\alpha\abs{\nabla_yG_L(t,J,y,0)}]^{2q}\right)^{1/(2q)}
 \le C_{\alpha,q}(1+t)^{-d/2-1/2+d/(4q)}\e^{-c t/L^2}.
 \label{eq:GNO-Theorem3b-specialized}
\end{equation}
The estimate is deterministic in $J$, uniform over all $L$-periodic coefficients in $[\lambda_*,1]$, and uniform in $L$.
\end{theorem}

The product torus law satisfies $\mathrm{SG}_L(\rho)$ by tensorization of its one-block spectral gap.  We now derive the lemma used in the main text.  Let $F_y(J)=F(\tau_yJ)$ and set $r=2q$, $r'=r/(r-1)$.  Choose $\alpha$ so that $\alpha r'>d$.  The environment--site intertwining and heat-kernel symmetry give, for each coordinate derivative,
\begin{equation}
 \abs{D\e^{-s\mathscr L_L}F}
 \le
 \left(\sum_y[\omega_{L,s}(y)^\alpha\abs{\nabla_yG_L(s,J,y,0)}]^r\right)^{1/r}
 \left(\sum_y\omega_{L,s}(y)^{-\alpha r'}\abs{F_y}^{r'}\right)^{1/r'}.
 \label{eq:weighted-Holder-full}
\end{equation}
By \cref{thm:GNO-Green-import}, the first factor is bounded by
\begin{equation}
 C(1+s)^{-d/2-1/2+d/(2r)}\e^{-cs/L^2}.
 \label{eq:first-Green-factor}
\end{equation}
For $2p\ge r'$, Minkowski and stationarity yield
\begin{align}
 &\left\|\left(\sum_y\omega_{L,s}(y)^{-\alpha r'}\abs{F_y}^{r'}\right)^{1/r'}\right\|_{L^{2p}(\Omega_L)}\\
 &\qquad\le
 \left(\sum_y\omega_{L,s}(y)^{-\alpha r'}\right)^{1/r'}\norm F_{L^{2p}(\Omega_L)}\\
 &\qquad\le C\min\{(1+s)^{d/(2r')},L^{d/r'}\}\norm F_{L^{2p}(\Omega_L)}.
 \label{eq:second-weight-factor}
\end{align}
The final sum estimate follows from $\alpha r'>d$ by comparison with the corresponding radial integral on the diffusive ball, capped by the torus volume.  Since $1/r+1/r'=1$, the powers cancel exactly:
\[
 -\frac d2-\frac12+\frac d{2r}+\frac d{2r'}=-\frac12.
\]
For $s\le L^2$ use the diffusive-volume term in \eqref{eq:second-weight-factor}; for $s\ge L^2$ use $L^{d/r'}\le(1+s)^{d/(2r')}$, while the exponential only improves the bound.  Consequently
\begin{equation}
 \norm{D\e^{-s\mathscr L_L}F}_{L^{2p}(\Omega_L)}
 \le C_p(1+s)^{-1/2}\norm F_{L^{2p}(\Omega_L)}.
 \label{eq:gradient-semigroup-upgrade}
\end{equation}
Apply this with $F=u(t/2)=\e^{-(t/2)\mathscr L_L}d_A$ and $s=t/2$.  Equation \eqref{eq:GNO-drift-moment} then gives
\begin{equation}
 \norm{D\e^{-t\mathscr L_L}d_A}_{L^{2p}}
 \le C_p\abs A(1+t)^{-1-d/4}.
\end{equation}
Under the exact stationary extension, this derivative is the physical oriented edge gradient $w_{x,\mu}(t)=b_{x,\mu}^*\e^{-tL_{J,L}}g_L$.  Spatial stationarity and restoration of $j_+$ prove the periodic gradient-propagator estimate in the main text for $p=1,2$.


\section{Reproducibility manifest}
\label{app:manifest}

The Paper~II directory contains the TeX and PDF, the seven fluctuation and sharp-cutoff verifiers, retained CSV/PNG/JSON certificates, the detailed gradient-propagator source map, a claim-status file, and SHA-256 checksums.  The root archive contains Paper~I and the reviewer-response matrix.

\begin{thebibliography}{99}

\bibitem{Sledgianowski2026RobustnessI}
N.~Sledgianowski,
\newblock Quenched pair-coupling disorder in frustration-free QGN parents I: exact weighted Hodge response, random-conductance homogenization, and annealed dynamical bounds,
\newblock companion research draft (2026), included in the consolidated archive.

\bibitem{GloriaMourrat2012}
A.~Gloria and J.-C.~Mourrat,
\newblock Spectral measure and approximation of homogenized coefficients,
\newblock \emph{Probability Theory and Related Fields} \textbf{154} (2012), 287--326.

\bibitem{GloriaNeukammOtto2013Long}
A.~Gloria, S.~Neukamm, and F.~Otto,
\newblock Quantification of ergodicity in stochastic homogenization: optimal bounds via spectral gap on Glauber dynamics---long version,
\newblock MiS Preprint 3/2013, Theorem 2 under (69b), Theorem 3(b), equation (74), and Corollary 8.

\bibitem{GloriaNeukammOtto2015}
A.~Gloria, S.~Neukamm, and F.~Otto,
\newblock Quantification of ergodicity in stochastic homogenization: optimal bounds via spectral gap on Glauber dynamics,
\newblock \emph{Inventiones Mathematicae} \textbf{199} (2015), 455--515.

\bibitem{deBuyerMourrat2015}
P.~de Buyer and J.-C.~Mourrat,
\newblock Diffusive decay of the environment viewed by the particle,
\newblock \emph{Electronic Communications in Probability} \textbf{20} (2015), 1--12.

\bibitem{RudelsonVershynin2013}
M.~Rudelson and R.~Vershynin,
\newblock Hanson--Wright inequality and sub-Gaussian concentration,
\newblock \emph{Electronic Communications in Probability} \textbf{18} (2013), no.~82, 1--9.

\bibitem{Dicke2021}
A.~Dicke,
\newblock Wegner estimate for random divergence-type operators monotone in the randomness,
\newblock \emph{Mathematical Physics, Analysis and Geometry} \textbf{24} (2021), article 22.

\bibitem{DickeVeselic2023}
A.~Dicke and I.~Veseli\'c,
\newblock Unique continuation for the gradient of eigenfunctions and Wegner estimates for random divergence-type operators,
\newblock \emph{Journal of Functional Analysis} \textbf{285} (2023), article 110040.

\bibitem{KamperEtAl2026}
M.~K\"amper, C.~Schumacher, F.~Schwarzenberger, and I.~Veseli\'c,
\newblock Quantitative concentration inequalities for the uniform approximation of the IDS,
\newblock arXiv:2602.18214 (2026).

\bibitem{Duerinckx2022Eigenvalue}
M.~Duerinckx,
\newblock Eigenvalue fluctuations for random elliptic operators in homogenization regime,
\newblock \emph{Journal of Statistical Physics} \textbf{187} (2022), article 32.

\bibitem{DuerinckxGloria2025}
M.~Duerinckx and A.~Gloria,
\newblock Large-scale dispersive estimates for acoustic operators: homogenization meets localization,
\newblock \emph{Journal of Functional Analysis} \textbf{289} (2025), article 111068.

\end{thebibliography}
\end{document}
